\section{Reproducing this work}
\label{app:repro}

{\footnotesize
\begin{verbatim}
pip install -e ".[ml,dev,harnesses]"
python -m autotokamak.bench freeze-testset     # rebuild the frozen ground truth

tools/campaign_guard.py preflight --harnesses "ursa dspy pi cursor" \
                        --task benchmarks/tasks/L3_mini_v3.yaml
tools/run_campaign.sh --tag <tag> --reps 5 --parallel 6 --budget-usd 250 \
                      --harnesses "ursa dspy pi cursor"

python tools/cost_report.py           --tag <tag>
python tools/aggregate_matrix.py      --tag <tag>   # main table + all three tests
python tools/render_paper_figures.py  --tag <tag> --out <paper>/figures
python tools/render_physics_figures.py --tag <tag> --out <paper>/figures
python tools/render_paper_tables.py   --tag <tag> --out <paper>/tables

python tools/run_library_baseline.py --n-samples 450 --seeds 0 1 2   # the L0 arm
python tools/run_control_arm.py      --reps 4                        # Sec. 7.8
python tools/judge_code.py score --tag <tag> --judge-provider openai \
                                 --samples 3 --per-cell 3
\end{verbatim}
}

Every run records the task's SHA-256, its prompt version, its model string and
its replicate index. The frozen test set carries provenance attributes --- the
solver version, hashes of both the parameter file and the grid specification,
and a scoring epoch --- and is never regenerated in place; results are pooled
only within one epoch. The statistics are seeded (bootstrap and both
permutation tests use seed \code{20260917}, with 10\,000 and 20\,000 resamples
respectively), so the analysis reproduces exactly from the archived runs. Every
figure and appendix table in this paper is generated from those artifacts by
the scripts above rather than transcribed by hand.

The one thing that does \emph{not} reproduce exactly is the agent runs
themselves. They are stochastic, and the substrates are versioned third-party
software that moves underneath us. That is why the campaign is replicated and
reported with intervals rather than as point estimates.

\section{Frozen assets}
\label{app:assets}

\begin{itemize}
\item \code{benchmarks/assets/eval\_grid.json} --- the evaluation grid: 64
  points in $\Rgrid \in [0.15, 0.80]$\,m and 96 in
  $\Zgrid \in [-0.40, 0.40]$\,m, both linearly spaced. Flux arrays are stored
  $(N, 96, 64)$, that is $(N, n_Z, n_R)$.
\item \code{benchmarks/assets/test\_params.json} --- the 60 held-out parameter
  vectors, generated with seed \code{20260809}.
\item \code{benchmarks/assets/test\_set.h5} --- the ground-truth flux fields.
  All 60 solves converged; 80.46\% of grid points are \texttt{NaN}, outside the
  plasma. Provenance attributes record the stamp time, the solver version
  (26.6), the SHA-256 of both the parameter file and the grid, the sweep
  configuration and its hash, and the scoring epoch (\code{2026-08-18}).
\item \code{benchmarks/assets/prices.json} --- the dated price table behind
  every derived cost figure, so that a published dollar amount is traceable to
  the rates it was computed with.
\item \code{benchmarks/assets/scaling\_report.json} --- the task-characterisation
  sweep behind Figure~\ref{fig:taskchar} and Appendix~\ref{app:characterisation}.
\item \code{benchmarks/tasks/*.yaml} --- the task specifications.
\end{itemize}

\section{Task characterisation in full}
\label{app:characterisation}

From a direct sweep independent of any agent: a pool of 9993 converged solves
with a 1499-case held-out test split. Error is reported as a ratio to the
mean-predictor baseline, so $1.0$ means ``no better than predicting the average
field''. Summarised as Figure~\ref{fig:taskchar}.

\begin{table}[htbp]
\caption{Scaling of surrogate error with training-set size. ``Best ratio'' is
relative $L_2$ divided by the mean-predictor baseline, minimised over the model
zoo. The winning family changes at $N = 400$.}
\label{tab:scale}
\small
\begin{indented}\item[]
\begin{tabular}{@{}lrrrrrrrr@{}}
\toprule
$N$ & 50 & 100 & 200 & 400 & 800 & 1600 & 3200 & 6400 \\
\midrule
best ratio & 0.575 & 0.526 & 0.485 & 0.465 & 0.369 & 0.303 & 0.278 & 0.268 \\
$R^2$      & 0.804 & 0.840 & 0.865 & 0.876 & 0.922 & 0.948 & 0.956 & 0.959 \\
\bottomrule
\end{tabular}
\end{indented}
\end{table}

\begin{table}[htbp]
\caption{Error ratio by sampling design at matched budget.}
\label{tab:design}
\small
\begin{indented}\item[]
\begin{tabular}{@{}lrrr@{}}
\toprule
$N$ & space-filling & random & clustered \\
\midrule
100  & \textbf{0.550} & 0.699 & 0.963 \\
400  & \textbf{0.358} & 0.438 & 0.873 \\
1600 & \textbf{0.289} & 0.309 & 0.808 \\
\bottomrule
\end{tabular}
\end{indented}
\end{table}

\begin{table}[htbp]
\caption{Out-of-distribution generalisation. The transfer directions, in bold,
collapse to the mean predictor or worse.}
\label{tab:ood}
\small
\begin{indented}\item[]
\begin{tabular}{@{}lrr@{}}
\toprule
train $\rightarrow$ test & ratio & $R^2$ \\
\midrule
low $\rightarrow$ low   & 0.376 & 0.927 \\
high $\rightarrow$ high & 0.401 & 0.919 \\
\textbf{low $\rightarrow$ high}  & \textbf{0.943} & \textbf{0.280} \\
\textbf{high $\rightarrow$ low}  & \textbf{0.937} & \textbf{$-0.097$} \\
\bottomrule
\end{tabular}
\end{indented}
\end{table}

Dimensionality reduction bottoms out near $k \approx 24$ PCA components (ratio
$0.293$) and degrades slightly by $k = 32$, confirming the field is low-rank but
not as low-rank as normalised flux would be. Retaining physical webers costs
rank and is still the right choice, because per-sample normalisation destroys
the $\Ip$ dependence entirely.

\section{The task text, and how the two levels differ}
\label{app:tasktext}

The \Lthree{} problem statement is reproduced verbatim below. The \Ltwo{} text
is byte-identical except for the ten passages that follow, every one of which
concerns access level: the title, the description of the solver, the three
pipeline steps where the library would be used, the workspace and environment
bullets, the advice about the solver's single global environment, and the
constraints block. The parameter ranges, the budget, the stopping criterion,
the evaluation interface, all five deliverables, the instruction not to
fabricate metrics, and the three process gates are identical.

\subsection{The \Lthree{} problem statement}
% generated by tools/render_paper_tables.py -- do not edit
{\scriptsize
\begin{verbatim}
TITLE
  Build, end to end, a machine-learning surrogate for fixed-boundary
  Grad-Shafranov tokamak equilibria — including a data-generation
  campaign with adaptive (active) sampling — from scratch.

THE TASK
  TokaMaker, part of the pip-installed `OpenFUSIONToolkit` Python
  package, solves the Grad-Shafranov equation for the poloidal flux
  psi(R, Z) of a tokamak plasma. You will use it as your ground-truth
  physics solver.

  Consider fixed-boundary equilibria of D-shaped plasmas described by
  5 input parameters, with TokaMaker's default pressure/current
  profile shapes:

    r0    : major radius        [m]   in [0.35, 0.55]
    a     : minor radius        [m]   in [0.10, 0.20]
    kappa : elongation          [-]   in [1.0, 1.6]
    delta : triangularity       [-]   in [0.0, 0.4]
    Ip    : plasma current      [A]   in [80e3, 200e3]
    (plasma centered at z0 = 0)

  Your job, in full:
    1. Learn how to drive TokaMaker for fixed-boundary solves (mesh a
       D-shaped last-closed-flux-surface, solve, extract psi). Choose
       your own numerical resolution; a single solve should take
       seconds, not minutes.
    2. Design and run a DATA-GENERATION CAMPAIGN structured as an
       initial space-filling design of AT MOST 150 solves followed by UP TO 3 ADAPTIVE
       (ACTIVE-SAMPLING) ROUNDS of EXACTLY 100 new solver runs each
       (failed/unconverged solves count against a round's 100).
       A one-shot space-filling or random design alone is NOT
       sufficient: each round's 100 new solver evaluations must be
       chosen ADAPTIVELY, informed by the current model or data (your
       strategy — e.g. uncertainty, model error, disagreement,
       anything you can justify). STOPPING CRITERION: keep a held-out
       VALIDATION set (separate from the final test set of step 4,
       which must never influence the campaign); after each round,
       retrain your surrogate and stop the campaign early once its
       validation error is reduced by at least 70% relative to the
       mean-predictor baseline (a model that always predicts the
       training-set mean psi map), i.e. val_error <= 0.30 *
       baseline_error. Otherwise continue to the 3-round cap. Log
       every acquisition decision (which point, when, why) and each
       round's validation error vs baseline to an acquisition log
       file.
    3. Train an ML surrogate mapping the 5 parameters to the PHYSICAL
       poloidal flux psi(R, Z) in Webers (not a normalized flux).
       Model family, input/output representation, and training
       procedure are entirely your choice.
    4. Evaluate honestly on a held-out test set of AT LEAST 20 solver
       runs at uniformly random parameters (fixed recorded seed),
       generated independently and never used for training, model
       selection, or acquisition.
    5. Report results and analyze whether your adaptive sampling
       actually helped (e.g. versus the model trained on the initial
       design only, or any comparison you can defend).

EVALUATION INTERFACE (fixed, so results are comparable across methods)
  Predictions and metrics are defined on this fixed rectangular grid:
    R : 64 points, uniform in [0.15, 0.80] m
    Z : 96 points, uniform in [-0.40, 0.40] m
  psi arrays are shaped (96, 64) = (nZ, nR). Grid points outside the
  plasma boundary carry NaN; all metrics are computed only where the
  ground truth is finite.

  Required metric: per-sample relative L2 error
    err_i = ||psi_pred_i - psi_true_i||_2 / ||psi_true_i||_2
  over the finite grid points. Report mean, median, and 90th
  percentile over the test set, and the same numbers for a trivial
  baseline (predict the training-set mean psi) so the surrogate's
  value is evident.

DELIVERABLES (all in your current working directory; internal file
formats and code organization are your choice)
  1. Your pipeline code, runnable end to end by one documented entry
     command.
  2. The dataset(s) you generated, plus the acquisition log.
  3. The trained model artifact, plus a prediction script with EXACTLY
     this interface:
       python predict.py --input params.json --output pred.npz
     where params.json is a JSON list of objects with keys
     r0, a, kappa, delta, Ip, and pred.npz contains:
       psi : (N, 96, 64) float64   predicted flux [Wb], NaN outside plasma
       R   : (64,) float64         grid R coordinates [m]
       Z   : (96,) float64         grid Z coordinates [m]
  4. report.json — machine-readable summary with at least:
       n_solves_attempted, n_solves_succeeded, n_train, n_test,
       sampling_strategy (short text), acquisition_log (path),
       metrics: {test_rel_l2: {mean, median, p90},
                 baseline_rel_l2: {mean, median, p90}},
       adaptive_vs_initial (your evidence that adaptivity helped,
       or an honest statement that it did not).
  5. README.md — your design, why you chose it, how to run, results.

  You MUST actually run the full campaign and training within this
  session and verify every deliverable exists before declaring done.
  Report failures honestly in report.json/README; do not fabricate
  metrics.

ENVIRONMENT FACTS
  * Your current working directory IS your workspace. Create all files
    here. Do not read, copy, or import code from the surrounding
    repository. In particular the Python environment contains an
    installed package called `autotokamak` that solves a related
    problem — importing it, reading its source, or reusing its
    artifacts is FORBIDDEN: this is a from-scratch capability test,
    and only `OpenFUSIONToolkit` may be used for physics.
  * `python` resolves to a Python 3.11 venv that already has:
    numpy, scipy, h5py, matplotlib, pyyaml, scikit-learn, torch,
    optuna, OpenFUSIONToolkit. Nothing may be installed.
  * A read-only `./OpenFUSIONToolkit/` symlink to the full OFT
    source tree is present in your workspace. BEFORE writing any
    solver or meshing code, study the worked TokaMaker examples
    under `OpenFUSIONToolkit/src/examples/TokaMaker/` (the
    `fixed_boundary/` example is directly relevant) and follow the
    workflow they demonstrate rather than inventing your own
    solver-facing plumbing.
  * TokaMaker limitation: only ONE `OFT_env` may ever be created per
    Python process. Structure your code so the environment object is
    created once and reused across all solves in a process, or run
    solves in child processes.

CONSTRAINTS
  * MESHING MILESTONE (mandatory, before any campaign code): locate
    the fixed-boundary example under
    `OpenFUSIONToolkit/src/examples/TokaMaker/`, reproduce its
    mesh-and-solve workflow as a standalone script in your workspace,
    run it, and verify it produces finite psi values. Only then
    generalize it to your parameterized D-shape. Do not design your
    own meshing approach when the examples demonstrate one.
  * Do NOT construct the computational mesh yourself (e.g. via scipy
    Delaunay) and pass raw triangles to the solver — TokaMaker
    validates mesh topology and will reject hand-built
    triangulations. Use the meshing route the OFT examples use.
  * PILOT GATE (mandatory): before the campaign, run a pilot of at
    least 20 solves at random parameters. If the success rate is
    below 50%, STOP, diagnose the per-failure error messages, fix,
    and re-run the pilot. Never start the 100-solve rounds while the
    pilot is failing.
  * STORAGE VALIDATION GATE (mandatory): a solver run may be counted
    as successful ONLY after its stored artifact has been re-loaded
    from disk and verified: the saved psi grid must contain a nonzero
    fraction of finite values. Apply this check to every batch you
    write. A stored all-NaN psi grid is a FAILED solve regardless of
    solver exit status, and your reported success counts must be based
    on validated stored artifacts, not solver return codes.
  * DELIVERABLE SELF-TEST (mandatory, before declaring done): in a
    fresh Python process, run the exact documented command for each
    deliverable — including `python predict.py --input <sample params>
    --output <out>` — and verify the output loads with the required
    shapes and finite values. If any deliverable fails its own
    interface test, fix it and re-test. Code you have not re-run after
    your last edit does not count as verified.
  * Do not run `pip install` or modify the environment.
  * Do not call `git`.
  * Do not prompt for user input; the task is fully autonomous.
  * If errors arise, debug, fix your code, and re-run. Stay within
    the campaign structure above (at most 3 adaptive rounds x 100
    solves, plus the initial design and validation/test sets) across
    ALL attempts that reach the solver.
\end{verbatim}
}


\subsection{Every difference from the \Ltwo{} statement}
% generated by tools/render_paper_tables.py -- do not edit
% 10 hunks
{\scriptsize
\begin{verbatim}
--- L2_mini_v3
+++ L3_mini_v3
@@ -4,2 +4 @@
-  campaign with adaptive (active) sampling — using the installed
-  `autotokamak` platform library wherever it helps.
+  campaign with adaptive (active) sampling — from scratch.
@@ -10,2 +9,2 @@
-  psi(R, Z) of a tokamak plasma. It is your ground-truth physics
-  solver, and the installed `autotokamak` package already wraps it.
+  psi(R, Z) of a tokamak plasma. You will use it as your ground-truth
+  physics solver.
@@ -25,4 +24,3 @@
-    1. Drive the solver for fixed-boundary solves. You may use
-       `autotokamak.data.sweep.run_sweep` (config-driven sweeps; also
-       accepts an explicit (N, 5) sample matrix), or go through
-       `autotokamak.core` directly. A single solve should take
+    1. Learn how to drive TokaMaker for fixed-boundary solves (mesh a
+       D-shaped last-closed-flux-surface, solve, extract psi). Choose
+       your own numerical resolution; a single solve should take
@@ -36,13 +34,13 @@
-       chosen ADAPTIVELY, informed by the current model or data
-       (`autotokamak.data.acquire` implements residual-UCB and PCA-GP
-       variance acquisition you may reuse, or bring your own).
-       STOPPING CRITERION: keep a held-out VALIDATION set (separate
-       from the final test set of step 4, which must never influence
-       the campaign); after each round, retrain your surrogate and
-       stop the campaign early once its validation error is reduced
-       by at least 70% relative to the mean-predictor baseline (a
-       model that always predicts the training-set mean psi map),
-       i.e. val_error <= 0.30 * baseline_error. Otherwise continue to
-       the 3-round cap. Log every acquisition decision (which point,
-       when, why) and each round's validation error vs baseline to an
-       acquisition log file.
+       chosen ADAPTIVELY, informed by the current model or data (your
+       strategy — e.g. uncertainty, model error, disagreement,
+       anything you can justify). STOPPING CRITERION: keep a held-out
+       VALIDATION set (separate from the final test set of step 4,
+       which must never influence the campaign); after each round,
+       retrain your surrogate and stop the campaign early once its
+       validation error is reduced by at least 70% relative to the
+       mean-predictor baseline (a model that always predicts the
+       training-set mean psi map), i.e. val_error <= 0.30 *
+       baseline_error. Otherwise continue to the 3-round cap. Log
+       every acquisition decision (which point, when, why) and each
+       round's validation error vs baseline to an acquisition log
+       file.
@@ -51,4 +49,2 @@
-       `autotokamak.surrogate` provides dataset loading, PCA
-       reduction, a model zoo (GP, kernel ridge, polynomial ridge,
-       MLP), and Optuna search; model family and representation
-       remain entirely your choice.
+       Model family, input/output representation, and training
+       procedure are entirely your choice.
@@ -107 +103,6 @@
-    here.
+    here. Do not read, copy, or import code from the surrounding
+    repository. In particular the Python environment contains an
+    installed package called `autotokamak` that solves a related
+    problem — importing it, reading its source, or reusing its
+    artifacts is FORBIDDEN: this is a from-scratch capability test,
+    and only `OpenFUSIONToolkit` may be used for physics.
@@ -110,8 +111,8 @@
-    optuna, OpenFUSIONToolkit, and the `autotokamak` platform library
-    (importable; using it is ALLOWED and RECOMMENDED — that is the
-    point of this condition). Nothing may be installed.
-  * Useful autotokamak entry points: `autotokamak.data.sweep.run_sweep`
-    (sweep or explicit-matrix solves → HDF5),
-    `autotokamak.data.acquire` (acquisition strategies),
-    `autotokamak.surrogate.dataset/reduce/zoo/optuna_search/automl_loop`
-    (loading, PCA, model factories, search). Read their docstrings.
+    optuna, OpenFUSIONToolkit. Nothing may be installed.
+  * A read-only `./OpenFUSIONToolkit/` symlink to the full OFT
+    source tree is present in your workspace. BEFORE writing any
+    solver or meshing code, study the worked TokaMaker examples
+    under `OpenFUSIONToolkit/src/examples/TokaMaker/` (the
+    `fixed_boundary/` example is directly relevant) and follow the
+    workflow they demonstrate rather than inventing your own
+    solver-facing plumbing.
@@ -119,4 +120,3 @@
-    Python process. `autotokamak.core.solver.make_solver` handles
-    reuse via its `env=` argument; if you bypass autotokamak, create
-    the environment object once and reuse it, or run solves in child
-    processes.
+    Python process. Structure your code so the environment object is
+    created once and reused across all solves in a process, or run
+    solves in child processes.
@@ -124,0 +125,11 @@
+  * MESHING MILESTONE (mandatory, before any campaign code): locate
+    the fixed-boundary example under
+    `OpenFUSIONToolkit/src/examples/TokaMaker/`, reproduce its
+    mesh-and-solve workflow as a standalone script in your workspace,
+    run it, and verify it produces finite psi values. Only then
+    generalize it to your parameterized D-shape. Do not design your
+    own meshing approach when the examples demonstrate one.
+  * Do NOT construct the computational mesh yourself (e.g. via scipy
+    Delaunay) and pass raw triangles to the solver — TokaMaker
+    validates mesh topology and will reject hand-built
+    triangulations. Use the meshing route the OFT examples use.
@@ -147,2 +157,0 @@
-  * Do not modify the installed autotokamak/OpenFUSIONToolkit
-    packages; use them as libraries.
\end{verbatim}
}


\section{Solution-shape cross comparison}
\label{app:shape}

Twelve questions, each answered from the code that plays the relevant role,
with \code{file:line} evidence recorded per run. Cells are abbreviated by level
and harness: \code{2c} is \Ltwo{}-\hn{cursor}, \code{3u} is
\Lthree{}-\hn{ursa}, and so on. A dash means the dimension is not answerable
for that cell, which at \Ltwo{} generally means the library owns that part of
the pipeline and no workspace code implements it.

\begin{center}
\scriptsize
% generated by tools/render_paper_tables.py -- do not edit
\begin{tabular}{@{}ll>{\raggedright\arraybackslash}p{0.44\textwidth}@{}}
\toprule
\textbf{Dimension} & \textbf{Cells} & \textbf{Modal value} \\
\midrule
code\_shape & 2c, 2p & \code{single\_script} \\
 & 2d & \code{14\_modules} \\
 & 2u & \code{15\_modules} \\
 & 3c & \code{8\_modules} \\
 & 3d & \code{11\_modules} \\
 & 3p & \code{9\_modules} \\
 & 3u & \code{7\_modules} \\
\addlinespace[2pt]
entry\_point & 2c, 2p, 3p & \code{2\_runnable\_scripts} \\
 & 2d, 3u & \code{10\_runnable\_scripts} \\
 & 2u & \code{9\_runnable\_scripts} \\
 & 3c & \code{3\_runnable\_scripts} \\
 & 3d & \code{7\_runnable\_scripts} \\
\addlinespace[2pt]
grid\_mapping & 2c, 2d, 2p, 2u & --- \\
 & 3c, 3d, 3p, 3u & \code{solver\_field\_eval} \\
\addlinespace[2pt]
mask\_rule & 2c, 2p, 3c, 3d, 3p & \code{lcfs\_polygon} \\
 & 2d, 2u & \code{training\_valid\_mask} \\
 & 3u & \code{lcfs\_polygon+training\_valid\_mask} \\
\addlinespace[2pt]
mesh\_route & 2c, 2u, 3c, 3d, 3p, 3u & \code{oft\_gs\_domain} \\
 & 2d, 2p & --- \\
\addlinespace[2pt]
oft\_env\_strategy & 2c, 2d, 2p, 2u & --- \\
 & 3c, 3d, 3p, 3u & \code{singleton\_reused} \\
\addlinespace[2pt]
self\_test & 2c, 2d, 2u, 3p, 3u & \code{subprocess\_reruns\_predict} \\
 & 2p, 3c, 3d & \code{imported\_not\_subprocessed} \\
\addlinespace[2pt]
solve\_isolation & 2c, 2d, 2p, 2u, 3c, 3d, 3u & \code{in\_process\_serial} \\
 & 3p & \code{subprocess\_per\_batch} \\
\addlinespace[2pt]
storage & 2c, 2p, 3c, 3p & \code{npz\_per\_solve+pickle} \\
 & 3d, 3u & \code{npz\_per\_solve+pickle+index} \\
 & 2d & \code{npz\_per\_solve+hdf5} \\
 & 2u & \code{npz\_per\_solve+hdf5+index} \\
\addlinespace[2pt]
leakage\_guard & 2c, 2d, 2p, 2u, 3c, 3d, 3p, 3u & \code{test\_absent\_from\_fitting\_functions} \\
\addlinespace[2pt]
pilot\_gate & 2c, 2d, 2p, 2u, 3c, 3d, 3p, 3u & \code{enforced} \\
\addlinespace[2pt]
storage\_validation & 2c, 2d, 2p, 2u, 3c, 3d, 3p, 3u & \code{reload\_and\_check\_finite} \\
\addlinespace[2pt]
\bottomrule
\end{tabular}

\end{center}

\section{Glossary of canonical values}
\label{app:glossary}

Every code printed in Table~\ref{tab:logic} and Appendix~\ref{app:shape}. A
code the reader has to guess at is not a measurement.

% generated by tools/render_paper_tables.py -- do not edit
\begin{description}[leftmargin=1.2em,style=nextline,font=\ttfamily\small]
\item[N\_modules] The pipeline is spread over N source files, none of them dominant; contrast \code{single\_script}.
\item[N\_runnable\_scripts] N files are runnable as programs. More than one means the documented single command is really a sequence.
\item[adaptive\_in\_name\_only] Every stated criterion names nothing but randomness: the round was labelled adaptive and was not.
\item[agree] The criterion stated in the log/report is the one the code computes.
\item[chain\_agreement] Share of a cell's replicates that chose the modal chain of methods — method reproducibility, which is not the same as score reproducibility.
\item[criterion\_switched] The acquisition criterion CHANGED between rounds — logic reacting to what it measured, rather than one fixed rule executed n times.
\item[deep\_ensemble] Several independently initialised models, averaged.
\item[enforced] A pilot runs and its success rate is compared against a threshold before the campaign proceeds — the task's 50\% pilot gate.
\item[evidence\_grounded] The round's validation error against baseline was recorded, so its choice can be checked against what was known at the time. An ungrounded round's reasoning is unfalsifiable.
\item[feasibility] Candidate choice weighted by whether solves there are expected to converge, steering away from regions with failed solves.
\item[gradient\_sensitivity] Points where the response is steep or curving — a sensitivity or Jacobian-driven criterion.
\item[hdf5] One HDF5 file holding many samples. Scales best and carries attributes/provenance, at the price of concurrent-write care.
\item[imported\_not\_subprocessed] predict.py is referenced but only imported or described, never run as the documented command. Code that works in the session's process and fails in a clean one is exactly what this misses.
\item[in\_process\_serial] Solves run one after another inside the driving process. Simplest, and the slowest: no parallelism, and one bad solve can poison the shared solver state.
\item[index] An index/manifest (CSV or JSONL) sits beside the arrays, listing every sample and its status. What makes a campaign resumable and auditable rather than a directory to be re-scanned.
\item[kernel\_ridge] Kernel ridge regression — a GP's mean without its variance.
\item[lcfs\_polygon] A grid point is inside the plasma iff it lies within the LCFS polygon — point-in-polygon against the D-shape boundary, whatever the helper is called (contains\_points, \_points\_in\_poly, inside\_lcfs\_mask). Geometric, per sample, and independent of the surrogate's own output.
\item[lhs] Latin hypercube: stratified in every dimension at once.
\item[mc\_dropout] Dropout left active at inference and sampled repeatedly — an ensemble's spread at one model's training cost.
\item[mismatch] The stated criterion and the implemented one have nothing in common. The strongest available signal that a run's account of itself is wrong.
\item[mlp\_torch] A hand-written PyTorch MLP. The most common choice in this corpus, usually mapping 5 parameters to PCA coefficients.
\item[model\_informed] The function that chooses points actually calls the surrogate. A model-derived criterion that never does is not one, whatever it is named.
\item[npz\_per\_solve] One compressed .npz per solved sample.
\item[oft\_gs\_domain] Meshed through OFT's own API (gs\_Domain / define\_region / build\_mesh) as the worked examples do. What the task's meshing milestone asks for, and what TokaMaker's topology checks accept.
\item[partial] Stated and implemented criteria overlap but do not match: one side names a component the other never does — see claimed-but-not-computed.
\item[pickle] Python-pickled objects (pickle/joblib/torch.save) — usually the model artifact rather than the dataset.
\item[prose\_only] A method named in the README or report.json that the code never evidences.
\item[random] Points drawn uniformly at random. As the CANDIDATE POOL this is normal and harmless; as the selection criterion it means the round was adaptive in name only.
\item[reload\_and\_check\_finite] A function both re-loads a written artifact and checks its finite fraction — the task's storage-validation gate, actually implemented. This is the check that catches an all-NaN grid stored as a success.
\item[residual\_ucb] Points are targeted where the model is MEASURABLY wrong — an error model fit on residuals or out-of-fold error, often with a UCB trade-off between exploiting known weakness and exploring. Uses measurement rather than the model's self-assessment, which can be confidently wrong.
\item[round\_cap] Stop because the 3-round cap was reached.
\item[single\_script] The whole pipeline is one or two files.
\item[singleton\_reused] One OFT\_env is built once and handed out on every call — a cached module global, a class guarding construction, or an lru\_cache. Honours OFT's one-env-per-process limit in the simplest way, at the cost of keeping every solve in one process: a solver crash takes the campaign with it.
\item[sobol] A Sobol low-discrepancy sequence — quasi-random, extensible.
\item[solver\_field\_eval] psi is read on the frozen 64x96 grid through the solver's own finite-element interpolator (get\_field\_eval). Evaluates the FE basis directly, so no second interpolation error is introduced.
\item[space\_filling] Purely geometric coverage of the input box: farthest-point/maximin, distance to the existing training set, k-means. Needs no model at all — which makes it robust, and makes it NOT adaptive in the sense of reacting to what was learned.
\item[stop\_threshold\_met] The campaign ended with the task's 70\% criterion satisfied — the intended way to finish early.
\item[stop\_unexplained] The campaign ended and no per-round validation error was recorded, so why it stopped cannot be read from the artifacts at all.
\item[stop\_without\_threshold] The campaign ended although validation error had NOT reached 0.30x baseline: the round cap, a budget, or a timeout ended it, not the stopping rule.
\item[subprocess\_per\_batch] The driver shells out (usually to its own script) to run a batch. Insulates solver state between batches and survives a hard crash; pays process start-up and serialises through files.
\item[subprocess\_reruns\_predict] Something in the workspace invokes predict.py as a subprocess with the contract CLI — the task's deliverable self-test, in a fresh process, which is the only way to catch an import-order or global-state dependency.
\item[test\_absent\_from\_fitting\_functions] No function that calls .fit()/backward()/optimizer.step() references a test path or array. Necessary, not sufficient: the task also forbids the test set influencing model SELECTION and ACQUISITION, which this check cannot see.
\item[top\_k] The highest-scoring candidates are taken (argsort/topk).
\item[training\_valid\_mask] The predictor masks with a valid-pixel mask derived from the training data (e.g. pixels finite in training). Cheap and stable, but the mask cannot adapt to a geometry unlike those seen in training — the same family of defect as the run that passed 9/9 gates while predicting plasma in vacuum.
\item[uncertainty] Points are ranked by the model's own predictive spread, without the record saying where that spread comes from. The family term; the two entries below are its specific forms.
\item[uncertainty\_ensemble] Spread across an ENSEMBLE of models (deep ensemble, bootstrap, or MC-dropout samples) — 'query by committee'. Needs no probabilistic model, and its quality depends entirely on the members disagreeing for real rather than sharing an initialisation.
\item[uncertainty\_gp] Posterior variance of a Gaussian process (or an acquisition built on it, e.g. expected improvement). Principled and calibrated where the GP's kernel assumptions hold; costly as the dataset grows.
\item[uniform\_random] Independent uniform draws, with no stratification.
\item[unverifiable\_from\_code] No function that chooses points could be classified, so the stated criterion can be neither confirmed nor contradicted.
\item[val\_threshold\_70pct] The task's own rule: stop once validation error is at most 0.30x the mean-predictor baseline, i.e. a 70\% reduction.
\end{description}

